% =====================================================================
% Preamble for eulerian_formal.tex
% A Stanley-style mathematical chapter with formal Rocq inline.
% =====================================================================

\documentclass[11pt,a4paper]{report}

% --- Encoding & fonts ---
\usepackage[T1]{fontenc}
\usepackage[utf8]{inputenc}
\usepackage{lmodern}
\usepackage{microtype}

% --- Math ---
\usepackage{amsmath, amssymb, amsthm, mathtools}
\usepackage{stmaryrd}  % \llbracket \rrbracket

% --- Page geometry ---
\usepackage[margin=1in]{geometry}

% --- Hyperlinks ---
\usepackage[colorlinks=true,
            linkcolor=blue!50!black,
            citecolor=blue!50!black,
            urlcolor=blue!50!black]{hyperref}
\usepackage{bookmark}

% --- Tables ---
\usepackage{booktabs}

% --- Code listings ---
\usepackage{xcolor}
\definecolor{rocqkw}{RGB}{120,30,150}
\definecolor{rocqcom}{RGB}{0,128,0}
\definecolor{rocqstr}{RGB}{200,30,30}
\definecolor{rocqbg}{RGB}{248,248,248}
\definecolor{rocqframe}{RGB}{220,220,230}
\definecolor{decodedbg}{RGB}{255,253,238}
\definecolor{sidebarbg}{RGB}{240,245,255}
\definecolor{sidebarframe}{RGB}{180,200,230}
\definecolor{computebg}{RGB}{245,255,245}

\usepackage{listings}

\lstdefinelanguage{Rocq}{
  morekeywords={Definition, Theorem, Lemma, Corollary, Proposition,
    Proof, Qed, Defined, Admitted, Fixpoint, Inductive, Variable,
    Variables, Hypothesis, Section, End, Module, Notation, Implicit,
    Types, Set, Unset, Arguments, From, Require, Import, Export,
    forall, exists, fun, match, with, end, if, then, else, let, in,
    return, as, structural, fix, cofix, struct, Type, Prop, Set,
    bool, nat, true, false, Some, None, option, list, prod, sum,
    Compute, Print, Check, Eval, Goal, Show, About, Reset, Open,
    Scope, Local},
  morekeywords=[2]{descent_set, des, asc, eulerian, beta, is_descent,
    descent, perm, set, ord_max, ord0, widen_ord, lift, val,
    insert_max_perm, insert_max_fun, extract_max_perm, ssreflect,
    fingroup},
  sensitive=true,
  morecomment=[s]{(*}{*)},
  morestring=[b]",
  basicstyle=\ttfamily\small,
  keywordstyle=\color{rocqkw}\bfseries,
  keywordstyle=[2]\color{rocqkw}\bfseries,
  commentstyle=\color{rocqcom}\itshape,
  stringstyle=\color{rocqstr},
  backgroundcolor=\color{rocqbg},
  frame=single,
  framerule=0.5pt,
  rulecolor=\color{rocqframe},
  framesep=4pt,
  xleftmargin=10pt,
  xrightmargin=10pt,
  breaklines=true,
  breakatwhitespace=true,
  showstringspaces=false,
  % UTF-8 chars that genuinely appear in our .v files (e.g. sigma in
  % qfact.v / qeul.v).  Without these, listings' default 8-bit handler
  % rejects the bytes with "Invalid UTF-8" errors.
  extendedchars=true,
  inputencoding=utf8,
  literate={->}{{$\to$}}{1}
           {=>}{{$\Rightarrow$}}{1}
           {>=}{{$\geq$}}{1}
           {<=}{{$\leq$}}{1}
           {<>}{{$\neq$}}{1}
           {forall}{{$\forall$}}{1}
           {exists}{{$\exists$}}{1}
           {σ}{{$\sigma$}}{1}
           {τ}{{$\tau$}}{1}
           {α}{{$\alpha$}}{1}
           {β}{{$\beta$}}{1}
           {γ}{{$\gamma$}}{1}
           {δ}{{$\delta$}}{1}
           {ε}{{$\varepsilon$}}{1}
           {π}{{$\pi$}}{1}
           {ω}{{$\omega$}}{1}
           {Σ}{{$\Sigma$}}{1}
           {ψ}{{$\psi$}}{1}
           {φ}{{$\phi$}}{1}
}

% --- (no lstnewenvironment wrapper; rocqbox uses lstlisting directly) ---

% --- Math notation aligned with Stanley ---
\newcommand{\Sn}[1]{\mathfrak{S}_{#1}}      % Stanley's S_n
\newcommand{\Snp}[1]{\mathfrak{S}_{#1}}     % alias
\newcommand{\des}{\mathrm{des}}
\newcommand{\asc}{\mathrm{asc}}
\newcommand{\maj}{\mathrm{maj}}
\newcommand{\inv}{\mathrm{inv}}
\newcommand{\descent}{\mathrm{D}}
\newcommand{\ord}[1]{\mathtt{'I_{#1}}}      % 'I_n
\newcommand{\Set}[1]{\mathtt{\{set\ #1\}}}
\newcommand{\Perm}[1]{\mathtt{\{perm\ #1\}}}
\newcommand{\eulerianA}{A}                  % Stanley's Eulerian A(n,k)
\newcommand{\Eul}{E}                        % Euler numbers
\newcommand{\beps}[2]{\beta_{#1}(#2)}       % Stanley's β_n(S)
\newcommand{\stirling}[2]{c(#1, #2)}        % Stanley's c(n,k)

% --- Theorem environments (Stanley-like) ---
\theoremstyle{plain}
\newtheorem{theorem}{Theorem}[chapter]
\newtheorem{proposition}[theorem]{Proposition}
\newtheorem{lemma}[theorem]{Lemma}
\newtheorem{corollary}[theorem]{Corollary}

\theoremstyle{definition}
\newtheorem{definition}[theorem]{Definition}
\newtheorem{example}[theorem]{Example}
\newtheorem{remark}[theorem]{Remark}

% --- Custom: formal/informal pairing ---
%
% \rocqsource{file}{line}  — clickable [Rocq, file:line] link to GitHub.
% Uses \GitHubBase and \GitHubBranch from below.

\newcommand{\GitHubBase}{https://github.com/LLM4Rocq/mathcomp-eulerian/blob}
\newcommand{\GitHubBranch}{main}

% \rocqsource takes a filename which may contain underscores.
% \detokenize keeps the underscores literal in BOTH the URL argument
% of \href AND the visible text (otherwise `_` is the subscript catcode
% in body text and triggers a math-mode error).
\newcommand{\rocqsource}[2]{%
  \href{\GitHubBase/\GitHubBranch/\detokenize{#1}\#L#2}%
       {\textsf{[\,Rocq, \texttt{\detokenize{#1}:#2}\,]}}%
}

\newcommand{\rocqsourcerange}[3]{%
  \href{\GitHubBase/\GitHubBranch/\detokenize{#1}\#L#2-L#3}%
       {\textsf{[\,Rocq, \texttt{\detokenize{#1}:#2-#3}\,]}}%
}

% Usage pattern:
%
%   \rocqsource{file.v}{LINE}
%   \begin{lstlisting}[language=Rocq]
%   ...formal Rocq text verbatim...
%   \end{lstlisting}
%
% (Wrapping listings in \newenvironment is fragile across multi-pass
%  builds with hyperref, so we use the listing environment directly.)

% --- Pedagogical sidebars ---

% \decoded{Plain-text mechanical translation of the formal block.}
% Appears immediately after a Rocq block; substitutes mathcomp casts for
% their meaning so the reader can see the formal text says what the
% informal text claims.
\newcommand{\decoded}[1]{%
  \par\nopagebreak\smallskip%
  \noindent\colorbox{decodedbg}{%
    \begin{minipage}{\dimexpr\textwidth-2\fboxsep\relax}%
      \footnotesize\textsf{Decoded.}\enspace #1%
    \end{minipage}%
  }\par\smallskip%
}

% \begin{sidebar}{Title} ... \end{sidebar}
% A "Why this representation?" or "Design choice" box, used to explain
% formalization decisions inline.  Implemented via a quote-style indent
% with a coloured rule on the left margin.
\newenvironment{sidebar}[1]{%
  \par\smallskip%
  \noindent\hspace{0.02\textwidth}%
  \begin{minipage}{0.96\textwidth}%
    \small%
    {\color{sidebarframe}\rule{2pt}{0.6\baselineskip}}\enspace%
    \textbf{Design choice: #1}\par\smallskip%
}{%
  \end{minipage}\par\smallskip%
}

% \begin{computestrip} ... \end{computestrip}
% A short Rocq Compute output strip, demonstrating a definition's value
% on a concrete example.
\lstnewenvironment{computestrip}{%
  \lstset{language=Rocq, backgroundcolor=\color{computebg}}%
}{}

% --- Title ---
\title{The mathcomp-eulerian formalization\\
  \large A Stanley-style reading of permutation statistics in Rocq/MathComp}
\author{}
\date{}
